\documentclass[11pt]{article}
\usepackage{amsmath,amssymb,amsthm,enumerate,nicefrac,fancyhdr,hyperref,graphicx,adjustbox,dsfont}
\hypersetup{colorlinks=true,urlcolor=blue,citecolor=blue,linkcolor=blue}
\usepackage[left=2.6cm, right=2.6cm, top=1.5cm, includehead, includefoot]{geometry}
\usepackage[dvipsnames]{xcolor}
\usepackage[d]{esvect}

%% commands
%% useful macros [add to them as needed]
% sets
\newcommand{\C}{{\mathbb{C}}} 
\newcommand{\N}{{\mathbb{N}}}
\newcommand{\Q}{{\mathbb{Q}}}
\newcommand{\R}{{\mathbb{R}}}
\newcommand{\Z}{{\mathbb{Z}}}
\newcommand{\F}{{\mathbb{F}}}

% bases
\newcommand{\mA}{\mathcal{A}}
\newcommand{\mB}{\mathcal{B}}
\newcommand{\mC}{\mathcal{C}}
\newcommand{\mD}{\mathcal{D}}
\newcommand{\mE}{\mathcal{E}}
\newcommand{\mL}{\mathcal{L}}
\newcommand{\mM}{\mathcal{M}}
\newcommand{\mO}{\mathcal{O}}
\newcommand{\mP}{\mathcal{P}}
\newcommand{\mS}{\mathcal{S}}
\newcommand{\mT}{\mathcal{T}}

% linear algebra
\newcommand{\diag}{\operatorname{diag}}
\newcommand{\adj}{\operatorname{adj}}
\newcommand{\rank}{\operatorname{rank}}
\newcommand{\spn}{\operatorname{Span}}
\newcommand{\proj}{\operatorname{proj}}
\newcommand{\prp}{\operatorname{perp}}
\newcommand{\refl}{\operatorname{refl}}
\newcommand{\tr}{\operatorname{tr}}
\newcommand{\nul}{\operatorname{Null}}
\newcommand{\nully}{\operatorname{nullity}}
\newcommand{\range}{\operatorname{Range}}
\renewcommand{\ker}{\operatorname{Ker}}
\newcommand{\col}{\operatorname{Col}}
\newcommand{\row}{\operatorname{Row}}
\newcommand{\cof}{\operatorname{cof}}
\newcommand{\Num}{\operatorname{Num}}
\newcommand{\Id}{\operatorname{Id}}
\newcommand{\ipb}{\langle \thinspace, \rangle}
\newcommand{\ip}[2]{\left\langle #1, #2\right\rangle} % inner products
\newcommand{\M}[2]{M_{#1\times #2}(\F)}
\newcommand{\RREF}{\operatorname{RREF}}
\newcommand{\cv}[1]{\begin{bmatrix} #1 \end{bmatrix}}
\newenvironment{amatrix}[1]{\left[\begin{array}{@{}*{\numexpr#1-1}{c}|c@{}}}{\end{array}\right]}
\newcommand{\am}[2]{\begin{amatrix}{#1} #2 \end{amatrix}}

% vectors
\newcommand{\vzero}{\vv{0}}
\newcommand{\va}{\vv{a}}
\newcommand{\vb}{\vv{b}}
\newcommand{\vc}{\vv{c}}
\newcommand{\vd}{\vv{d}}
\newcommand{\ve}{\vv{e}}
\newcommand{\vf}{\vv{f}}
\newcommand{\vg}{\vv{g}}
\newcommand{\vh}{\vv{h}}
\newcommand{\vl}{\vv{\ell}}
\newcommand{\vm}{\vv{m}}
\newcommand{\vn}{\vv{n}}
\newcommand{\vp}{\vv{p}}
\newcommand{\vq}{\vv{q}}
\newcommand{\vr}{\vv{r}}
\newcommand{\vs}{\vv{s}}
\newcommand{\vt}{\vv{t}}
\newcommand{\vu}{\vv{u}}
\newcommand{\vvv}{{\vv{v}}}
\newcommand{\vw}{\vv{w}}
\newcommand{\vx}{\vv{x}}
\newcommand{\vy}{\vv{y}}
\newcommand{\vz}{\vv{z}}

% display
\newcommand{\ds}{\displaystyle}
\newcommand{\qand}{\quad\text{and}}
\newcommand{\qandq}{\quad\text{and}\quad}
\newcommand{\hint}{\textbf{Hint: }}
\newcommand{\tri}{\triangle}
\newcommand{\dso}{\mathds{1}}

% misc
\newcommand{\area}{\operatorname{area}}
\newcommand{\vol}{\operatorname{vol}}
\newcommand{\red}[1]{{\color{red} #1}}
\newcommand{\rc}{\red{\checkmark}}

\title{STAT 333 Notes}
\author{Thomas Liu}
\begin{document}
\maketitle
\tableofcontents

\newpage
\section{Review of Elementary Probability}
\subsection{Probability Space}
\subsubsection*{Experiment}
A fair coin is tossed once 
\subsubsection*{Sample Space}
Set of all possible outcomes ($\Omega$), $\Omega\in\{H,T\}$
\subsubsection*{Events}
Roughly set of all subsets of sample space $\Omega$, $\epsilon = \{\emptyset, \{H\}, \{T\}, \{H,T\}\}$
\subsection{Probability}
Probability is a function from the set of all events $\epsilon$ to $\R$, $P:\epsilon\to\R$
\begin{itemize}
    \item for any $A\in\epsilon$, $P(A)\in[0,1]$
    \item $P(\Omega)=1$
    \item for any $n\in\Z^+$, $E_1,E_2,\cdots,E_n\in\epsilon$ such that $E_i\cap E_j=\emptyset$ for $i\neq j$, then $\ds P(\cup_{i=1}^n E_i)=\sum_{i=1}^n P(E_i)$
\end{itemize}
\subsubsection*{Properties}
\begin{itemize}
    \item $P(\emptyset)=0$
    \item $A,B\in\epsilon$, $P(A\cup B)=P(A)+P(B)-P(A\cap B)$
    \item for any $A\in\epsilon$, $P(A^c)=1-P(A)$
    \item if $A,B$ satisfies $A\cap B=\emptyset$, then $P(A\cup B)=P(A)+P(B)$
    \item If $A\subseteq B$, then $P(A)\leq P(B)$
\end{itemize}
\subsection{Conditional Probability}
The conditional probability of an event $A$ given $B$ is 
\[P(A|B)=\frac{P(A\cap B)}{P(B)}\]
and is only defined when $P(B)>0$
\subsubsection*{Properties}
\begin{itemize}
    \item $P(A\cap B) = P(A|B)P(B)$
    \item for any $A_1,A_2,\cdots,A_n$, $P(A_1\cap A_2\cap\cdots\cap A_n)=P(A_1)P(A_2|A_1)P(A_3|A_1\cap A_2)\cdots P(A_n|A_1\cap A_2\cap\cdots\cap A_{n-1})$
    \item $P(A^c|B)=1-P(A|B)$
    \item $A_1,A_2$ are two events, $P(A_1\cup A_2|B)=P(A_1|B)+P(A_2|B)-P(A_1\cap A_2|B)$
\end{itemize}
\subsection{Independent Events}
Two events $A$ and $B$ are mutually independent of each other if $P(A\cap B)=P(A)P(B)$ \\
The events $A_1,A_2,\cdots,A_n$ are mutually independent if $\ds P(A_1\cap A_2\cap\cdots\cap A_n)=\prod_{i=1}^{n}P(A_i)$ 
\begin{itemize}
    \item if $A$ and $B$ are independent of each other, $P(A|B) = P(A)$
\end{itemize}
\subsection{Law of Total Probability}
Say $E_1,\cdots,E_n$ are disjoint events. If satisfies $E_i\cap E_j=\emptyset$ for $i\neq j$, $\cup_{i=1}^n E_i=\Omega$. Then for any event $A$, 
\[P(A) = \sum_{i=1}^n P(A|E_i)P(E_i)\]
\subsection{Bayes' Formula}
\begin{align*}
    P(E_i|A) &= \frac{P(A\cap E_i)}{P(A)} \\
             &= \frac{P(A\cap E_i)}{\sum_{j=1}^n P(A\cap E_j)} \\
             &= \frac{P(A|E_i)P(E_i)}{\sum_{j=1}^{n}P(A|E_j)P(E_j)}
\end{align*}
\subsection{Random Variables}
A function from sample space to real numbers $X:\omega\to\R$. Range of $X$: state space 
\subsection{Discrete Random Variables}
If state space / range of random variable is discrete ($\{a_1,a_2,\cdots\}$)
\subsubsection*{Probability Mass Function (PMF)}
\[P(a) = P(X=a)\]
\subsubsection*{Cumulative Distribution Function (CDF)}
\[F(a) = P(X\leq a) = \sum_{x\leq a}P(X=x)\]
\subsubsection*{Tail Probability Function (TPF)}
\[\overline{F}(a) = P(X>a) = 1-F(a)\]
\subsubsection*{Properties}
State space = $\{a_1,a_2,\cdots\}$, $a_1<a_2<\cdots$
\[p(a_i) = F(a_i) - F(a_{i-1})\]
\subsection{Discrete Distributions}
\subsubsection*{Bernoulli Distribution}
$x\in\{0,1\}$, $X\sim Ber(p)$
\[p(x) = p^x(1-p)^{1-x}\]
\subsubsection*{Binomial Distribution}
$x$: number of successes in $n$ independent $Ber(p)$ trials, $x\in\{0,1,\cdots,n\}$, $X\sim Bin(n, p)$
\[p(x) = \binom{n}{x}p^x(1-p)^{n-x}\]
\subsubsection*{Geometric Distribution}
$x_1,x_2,\cdots,\overset{iid}{\sim}Ber(p)$ \\
$y$: number of trials required to get the first success, $y\in\{1,2,\cdots\}$, $Y\sim Geo(p)$
\[p_Y(x) = P(Y=x) = (1-p)^{x-1}p\]
$z$: number of failures until the first success, $z\in\{0,1,\cdots\}$, $Z\sim Geo(p)$
\[p_Z(x) = P(Z=x) = (1-p)^{x}p\]
\subsubsection*{Poisson Distribution}
$X\sim Poi(\lambda)$, $x\in\{0,1,\cdots\}$, $\lambda>0$
\[p(x) = \frac{\lambda^x}{x!}e^{-\lambda}\]
\subsection{Continuous Random Variables}
$x$ is continuous random variable, $x$ takes values over continuous in $\R$ 
\subsubsection*{Probability Density Function (PDF)}
$f_X:\R\to[0,\infty)$, $\ds\int_{-\infty}^{\infty}f_X(x)dx=1$, $\ds F(x)=P(X\leq x)=\int_{-\infty}^{x}f_X(t)dt$, $f_X(x)=\ds\frac{dF(x)}{dx}$
\begin{enumerate}
    \item For continuous $x$, cdf $F(x)$ is continuous 
    \item For any interval $B\subseteq\R$ \[P(x\in B)=\int_B f_X(x)dx\], if $B\in(-\infty,x]\rightarrow P(x\in B)=F(x)$
\end{enumerate}
\subsubsection*{Uniform Distribution}
$X\sim U(a,b)$, $a<b$
\begin{align*}
    f_X(x) &= \begin{cases}
        \ds\frac{1}{b-a} & a\leq x\leq b \\
        0 & \text{otherwise}
    \end{cases}
\end{align*}
\subsubsection*{Exponential Distribution}
$X\sim Exp(\lambda)$, $\lambda>0$
\begin{align*}
    f_X(x) &= \begin{cases}
        \ds\lambda e^{-\lambda x} & x\geq 0 \\
        0 & \text{otherwise}
    \end{cases}
\end{align*}
\subsubsection*{Erlang Distribution}
$X\sim Erlang(n,\lambda)$, $n\in\{1,2,\cdots\}$, $\lambda>0$
\begin{align*}
    f_X(x) &= \begin{cases}
        \ds\frac{\lambda^n}{(n-1)!}x^{n-1}e^{-\lambda x} & x\geq 0 \\
        0 & \text{otherwise}
    \end{cases}
\end{align*}
\subsubsection*{Gamma Distribution}
$X\sim Gamma(\alpha,\lambda)$, $\alpha>0$, $\lambda>0$
\begin{align*}
    f_X(x) &= \begin{cases}
        \ds\frac{\lambda^{\alpha}}{\Gamma(\alpha)}x^{\alpha-1}e^{-\lambda x} & x\geq 0 \\
        0 & \text{otherwise}
    \end{cases}
\end{align*}
\subsection{Expectation}
$X$ is a continuous random variable, $g:\R\to\R$ is a function 
\begin{align*}
    E[g(X)] &= \begin{cases}
        \ds\int_{-\infty}^{\infty}g(x)f_X(x)dx & \text{if } g(x) \text{ is continuous} \\
        \ds\sum_{x\in\R}g(x)p_X(x) & \text{if } g(x) \text{ is discrete}
    \end{cases}
\end{align*}
Some special $g(x)$
\begin{enumerate}
    \item $g(x)=x^n\Rightarrow E[g(X)] = E[X^n]$, $n^{th}$ moment of $X$ 
    \item $\mu=E(X)$, $g(x)=(x-\mu)^2\Rightarrow E[g(X)]=E[(X-\mu)^2]=Var(X)$
    \item $Var(X)=E(X^2)-[E(X)]^2$ 
    \item $E(aX+b)=aE(X)+b$, $a,b\in\R$
    \item $Var(aX+b)=a^2Var(X)$, $a,b\in\R$
\end{enumerate}
\subsection{Moment Generating Function (MGF)}
\subsubsection*{Properties}
\begin{itemize}
    \item $M_X(t)=E[e^{tX}]$, $t\in\R$
    \item $\dfrac{d^n}{dt}M_X(t)\bigg|_{t=0}=E[X^n]$ 
    \item $x,y$ are two random variables and identically distributed iff $M_X(t)=M_Y(t)$
    \item $x_1,x_2,\cdots,x_n$ are independent, then \[M_{X_1+X_2+\cdots+X_n}(t)=M_{X_1}(t)M_{X_2}(t)\cdots M_{X_n}(t) = \prod_{i=1}^{n}M_{X_i}(t)\]
\end{itemize}
\subsection{Joint Distribution}
Two random variables $X$ and $Y$ defined on the same sample space \\
cdf / Joint cdf: \[F_{X,Y}(x,y)=P(X\leq x, Y\leq y)\] 
Marginal cdf: \[F_X(x)=P(X\leq x)=\lim_{y\to\infty}F_{X,Y}(x,y),\ F_Y(y)=P(Y\leq y)=\lim_{x\to\infty}F_{X,Y}(x,y)\]
\subsubsection{Discrete Joint Distribution}
Joint pmf: $p_{X,Y}(x,y)=P(X=x, Y=y)$ 
\[P_X(x)=\sum_{y}p_{X,Y}(x,y),\ P_Y(y)=\sum_{x}p_{X,Y}(x,y)\] 
\subsubsection{Continuous Joint Distribution}
Joint pdf: $f_{X,Y}(x,y)=\dfrac{\partial^2}{\partial x \partial y}F_{X,Y}(x,y)$ 
\[F_{X,Y}(x,y)=\int_{-\infty}^{x}\int_{-\infty}^{y}f_{X,Y}(t,s)dtds\]
\subsection{Expectation}
\begin{align*}
    E[g(X,Y)] &= \begin{cases}
        \ds\int_{-\infty}^{\infty}\int_{-\infty}^{\infty}g(x,y)f_{X,Y}(x,y)dydx & \text{if } g(x,y) \text{ is continuous} \\
        \ds\sum_{x\in\R}\sum_{y\in\R}g(x,y)p_{X,Y}(x,y) & \text{if } g(x,y) \text{ is discrete}
    \end{cases}
\end{align*}
$E[g(X,Y)]$ exists if one of the following happens
\begin{itemize}
    \item $E[g(X,Y)]<\infty$
    \item $g(X,Y)\geq0$
    \item $g(X,Y)\leq0$
\end{itemize}
\subsection{Independence of Random Variables}
$X$ and $Y$ are independent if $F_{X,Y}(x,y)=F_X(x)F_Y(y)$ \\
For discrete case: \[P_{X,Y}(x,y)=P_X(x)P_Y(y)\] 
For continuous case: \[f_{X,Y}(x,y)=f_X(x)f_Y(y)\]
For expectation, if $X$ and $Y$ are independent, then \[E[g(X,Y)]=E[g(X)]E[g(Y)]\]
\subsection{Covariance}
\begin{align*}
    Cov(X,Y) &= E[(X-\mu_X)(Y-\mu_Y)] \\
             &= E[XY] - E[X]E[Y]
\end{align*}
\subsubsection{Properties}
\begin{itemize}
    \item If $X$ and $Y$ are independent, then $Cov(X,Y)=0$, but $Cov(X,Y)=0$ does not imply $X$ and $Y$ are independent
    \item $X$ and $Y$ are independent, then $M_{X,Y}(t,s)=M_X(t)M_Y(s)$
    \item $X$ and $Y$ are independent, then $Var(aX+bY)=a^2Var(X)+b^2Var(Y)$
\end{itemize}
\subsection{Convergence Modes}
$X_1,X_2,\cdots$ is a sequence of random variables, $X$ is another random variable 
\begin{enumerate}
    \item $X_n\to X$ almost surely $P(X_n\to X)=1$. $\omega\in\Omega$ (one element of sample space), $X_n(\omega)\to X(\omega)$ as $n\to\infty$. $E=\{\omega\in\Omega: X_n(\omega)\to X(\omega)\}$, $P(E)=1$
    \item $X_n\to X$ in probability if for any $\epsilon>0$, $P(|X_n-X|>\epsilon)\to 0$ as $n\to\infty$
    \item $X_n\to X$ in distribution if $F_{X_n}(a)\to F_X(a)$ for every $a$ where $F_X(a)$ is continuous
\end{enumerate}
\subsubsection{Property}
$X_n\to X$ almost surely $\Rightarrow$ $X_n\to X$ in probability $\Rightarrow$ $X_n\to X$ in distribution \\
\textbf{Not the other way around}
\subsection{Strong Law of Large Numbers (SLLN)}
$X_1,X_2,\cdots$ are a sequence of iid random variables with $E(|X_1|)<\infty$, then 
\[\dfrac{X_1+\cdots+X_n}{n}\to E(X_1)\text{ almost surely as }n\to\infty\]
\newpage 

\section{Conditional Distributions and Conditional Expectation}
\subsection{Conditional Probability}
\subsubsection*{Discrete Case}
$X_1,X_2$ are discrete random variables, the conditional pmf of $X_1$ given $X_2$ is
\[P_{X_1|X_2}(x_1|x_2)=P(X_1=x_1|X_2=x_2)=\frac{P(X_1=x_1,X_2=x_2)}{P(X_2=x_2)}\]
is defined only when $P(X_2=x_2)>0$ 
\subsubsection*{Properties}
If $P_{X_1}(x_1), P_{X_2}(x_2)>0$, 
\[P_{X_2|X_1}(x_2|x_1)=\frac{P_{X_1|X_2}(x_1|x_2)P_{X_2(x_2)}}{P(X_1=x_1)}\]
If $X_1,X_2,X_3$ are discrete and $P_{X_3}(x_3)>0$, then
\[P_{X_1,X_2|X_3}(x_1,x_2|x_3)=\frac{P_{X_1,X_2,X_3}(x_1,x_2,x_3)}{P_{X_3}(x_3)}\]
\subsubsection*{Continuous Case}
$X$ and $Y$ are two continuous random variables, 
\[f_{X,Y}(x,y): \text{joint pdf of }X,Y\]
\[f_X(x), f_Y(y): \text{marginal pdf of }X,Y\]
The conditional pdf of $X|Y=y$ is 
\[f_{X|Y}(x|y)=\frac{f_{X,Y}(x,y)}{f_Y(y)}\]
is defined only when $f_Y(y)>0$
\subsubsection*{Properties}
Assume that both $f_X(x)$ and $f_Y(y)$ are positive, then
\begin{enumerate}
    \item $\displaystyle f_{Y|X}(y|x)=\frac{f_{X|Y}(x|y)f_Y(y)}{f_X(x)}$
    \item $\displaystyle f_X(x)=\int f_{X|Y}(x|y)f_Y(y)dy$
    \item $\displaystyle f_Y(y)=\int f_{Y|X}(y|x)f_X(x)dx$
    \item $\displaystyle f_{Y|X}(y|x) = \frac{f_{X|Y}(x|y)f_Y(y)}{\int f_{X|Y}(x|y)f_Y(y)dy}$
\end{enumerate}
\subsection{Conditional Expectation}
\subsubsection*{Discrete Case}
\[E[X_1|X_2=x_2]=\sum_{x_1}x_1P_{X_1|X_2}(x_1|x_2)\]
More generally, if $w(\cdot,\cdot)$, $h(\cdot)$, and $g(\cdot)$ are functions, then
\[E(w(X_1,X_2)|X_2=x_2)=E(w(X_1,x_2)|X_2=x_2)=\sum_{x_1}w(x_1,x_2)P_{X_1|X_2}(x_1|x_2)\]
and 
\[E(g(X_1)h(X_2)|X_2=x_2)=E[g(X_1)h(x_2)|X_2=x_2]=h(x_2)E[g(X_1)|X_2=x_2]\]   
\[E[ag(X_1)+bh(X_1)|X_2=x_2]=aE[g(X_1)|X_2=x_2]+bE[h(X_1)|X_2=x_2]\]
\subsubsection*{Properties}
If $X_1,X_2,\cdots,X_n,Y$ jointly discrete random variables, and $P_Y(y)>0$, and for any $a_1,\cdots,a_n\in\R$, then
\[E[a_1X_1+a_2X_2+\cdots+a_nX_n|Y=y]=\sum_{i=1}^{n}a_iE[X_i|Y=y]\]
\subsubsection*{Continuous Case}
$E[g(Y)|X=x] = \ds\int g(y)f_{Y|X}(y|x)dy$ 
Special case: 
\begin{itemize}
    \item $g(y)=y$, then $E[Y|X=x]=\ds\int yf_{Y|X}(y|x)dy$
    \item $g(y)=(Y-E[Y|X=x])^2$, then $E[(Y-E[Y|X=x])^2|X=x]=Var(Y|X=x)$
\end{itemize}
\subsection{Independence}
If $X_1,X_2$ are independent, then 
\begin{itemize}
    \item $P_{X_1|X_2}(x_1|x_2)=P_{X_1}(x_1)$
    \item $E[X_1|X_2]=E[X_1]$
\end{itemize}
\subsection{Mix of Discrete and Continuous Random Variables}
$X$ is continuous random variable, and $Y$ is discrete random variable \\
Situation 1: $f_X(x)$ and $P_{Y|X}(y|x)$ are given, then 
\[f_{X|Y}(x|y) = \frac{f_{X,Y}(x,y)}{P_Y(y)}=\frac{f_X(x)P_{Y|X}(y|x)}{P_Y(y)}\]
Situation 2: $P_Y(y)$ and $f_{X|Y}(x|y)$ are given, then
\[P_{Y|X}(y|x)=\frac{f_{X|Y}(x|y)P_Y(y)}{f_X(x)}\]
\subsection{Compute Expectation by Conditioning}
$X$ and $Y$ are two random variables, and $g:\R\to\R$, then 
\[E[g(X)]=E[E[g(X)|Y]]\]
\subsubsection*{Thoerem}
If $X$ and $Y$ are two random variables, then 
\[Var(X)=E[Var(X|Y)]+Var(E[X|Y])\]
\subsection{Calculating Probabilities using Conditioning}
Suppose we have two continuous random variables $X$ and $Y$
\begin{align*}
    P(X<Y) &= \int P(X<Y|Y=y)f_Y(y)dy \\
           &= \int P(X<y|Y=y)f_Y(y)dy \\
           &\text{if $X$ and $Y$ are independent} \\
           &= \int F_X(y)f_Y(y)dy \\
           &= E[F_X(Y)]
\end{align*}
If $X$ and $Y$ are identically distributed, then
\begin{align*}
    P(X<Y) &= \int P(X<Y|Y=y)f_Y(y)dy \\
           &= \int P(Y>X|Y=y)f_Y(y)dy \\
           &= \int (1-F_X(y))f_Y(y)dy \\
           &= 1-E[F_X(Y)] \\
           &= 1/2
\end{align*}

\newpage 
\section{Discrete-Time Markov Chains}
\subsection{Definitions and Basic Concepts}
\subsubsection{Stochastic Process}
$\{X(t),t\in T\}$ is called a stochastic process if $X(t)$ is a rv for any given $t\in T$. $T$ is referred to as the index set and is often interpreted in the context of time 
\subsubsection{Discrete-Time Markov Chain (DTMC), Markov Property}
A stochastic process $\{X_n,n\in\N\}$ is said to be a discrete-time Markov chain if the following conditions hold true:
\begin{itemize}
    \item For $n\in\N$, $X_n$ is discrete rv 
    \item For $n\in\N$ and all states $x_0,x_1,\cdots,x_{n+1}\in S$, the Markov property must hold:
    \[P(X_{n+1}=x_{n+1}|X_n=x_n,X_{n-1}=x_{n-1},\cdots,X_0=x_0)=P(X_{n+1}=x_{n+1}|X_n=x_n)\]
    The current state $X_n$ contains all the information needed to predict the next state $X_{n+1}$, regardless of the previous states $X_{n-1},\cdots,X_0$
\end{itemize}
\subsubsection{Transition Probability}
For any pair of states $i$ and $j$, the transition probability from state $i$ at time $n$ to state $j$ at time $n+1$ is given by 
\[P_{n,i,j}=P(X_{n+1}=j|X_n=i)\]    
\subsubsection{Stationary Transition Probability}
A DTMC is called stationary or homogeneous if the transition probabilities do not depend on time 
\[P(X_n=j|X_{n-1} = i) = P(X_{n-1}=j|X_{n-2} = i) = \cdots = P(X_1=j|X_0=i) = P_{i,j}\]
\subsubsection{Transition Probability Matrix}
\subsubsection*{One Step Transition Probability Matrix}
Let $P_n$ be the associated matrix containing all these transition probabilities, referred to as the one-step transition probability matrix (TPM) from time $n$ to time $n+1$. It looks like 
\begin{align*}
    P_n = [P_{n,i,j}] = 
    \begin{bmatrix}
        P_{n,0,0} & P_{n,0,1} & \cdots & P_{n,0,j} \\
        P_{n,1,0} & P_{n,1,1} & \cdots & P_{n,1,j} \\
        \vdots & \vdots & \ddots & \vdots \\
        P_{n,i,0} & P_{n,i,1} & \cdots & P_{n,i,j} \\
        \vdots & \vdots & \ddots & \vdots 
    \end{bmatrix}
\end{align*}
\subsubsection*{Stationary and Homogeneous}
For each pair of states $i$ and $j$, if $P_{n,i,j}$ does not depend on $n$, $P_{n,i,j}=P_{i,j}$, $\forall n\in\N$, then we say the DTMC is stationary or homogeneous, one-step TPM is 
\begin{align*}
    P = [P_{i,j}] =
    \begin{bmatrix}
        P_{0,0} & P_{0,1} & \cdots & P_{0,j} \\
        P_{1,0} & P_{1,1} & \cdots & P_{1,j} \\
        \vdots & \vdots & \ddots & \vdots \\
        P_{i,0} & P_{i,1} & \cdots & P_{i,j} \\
        \vdots & \vdots & \ddots & \vdots
    \end{bmatrix}
\end{align*}
\begin{align*}
    P_{n,i,j}&=P(X_{n+1}=j|X_n=i)\\
    &= P(X_1=j|X_0=i) \\
    &= P_{i,j} 
\end{align*}
\subsubsection*{N-Step Transition Probability}
For any pair of states $i$ and $j$, the n-step transition probability is given by 
\[P_{i,j}^{(n)} = P(X_n = j | X_0 = i) = P(X_{m+n}=j|X_m=i)\quad m,n\in\N\]
Due to stationary assumption, this quantity is actually independent of $m$. Thus $P_{i,j}^{(n)} = P(X_n=j|X_0=i), n\in\N$ 
\begin{align*}
    P_{i,j}^{(n)} = P(X_0=j|X_0=i) = \begin{cases}
        0 & \text{if } i\neq j \\
        1 & \text{if } i=j
    \end{cases}
\end{align*} 
\subsubsection{Properties of Stationary Transition Matrix}
Say $P = [P_{i,j}]_{i,j\in S}$ is stationary transition matrix, then 
\begin{itemize}
    \item $P_{i,j}\geq 0$ 
    \item $\sum_j P_{i,j}=1$ 
\end{itemize}
\subsubsection{Chapman-Kolmogorov Equation}
For $n\in\Z^+$ 
\begin{align*}
    P_{ij}^{(n)} &= P(X_n=j|X_0=i) \\ 
                 &= \sum_{k=0}^\infty P(X_n=j|X_{n-1}=k,X_0=i) P(X_{n-1}=k|X_0=i) \\
                 &= \sum_{k=0}^{\infty} P_{i,k}^{(n-1)}P(X_n=j|X_{n-1}=k,X_0=i) \\
                 &= \sum_{k=0}^{\infty} P_{i,k}^{(n-1)}P(X_n=j|X_{n-1}=k) \\
                 &= \sum_{k=0}^{\infty} P_{i,k}^{(n-1)}P_{k,j}
\end{align*}
More generality, for $m,n\in\Z^+$,
\[P_{i,j}^{(m+n)} = \sum_{k\in S}P_{ik}^{(m)}P_{kj}^{(n)}\]
\subsubsection*{Property}
\[P^{(n)} = P\cdot P^{(n-1)} = P^n\]
\subsubsection{Initial Distribution \& n-Step Distribution}
$\alpha_0$: the vector consisting of all the pmf of $X_0$, the initial distribution of DTMC \\
$\alpha_n$: the vector consisting of all the pmf of $X_n$, the n-step distribution of DTMC 
\[\alpha_n = \alpha_0\cdot P^{(n)} = \alpha_0\cdot P^n\] 
\subsubsection{Accessible, Communicate}
State $j$ is accessible from state $i$ (denote by $i\to j$) if $\exists n\in\N$ such that 
$P_{i,j}^{(n)}>0$. If states $i$ and $j$ are accessible from each other, then states $i$ and $j$ communicate (denote by $i\leftrightarrow j$), 
$i\leftrightarrow j$ iff $\exists m,n\in\N$ such that $P_{i,j}^{(n)}>0$ and $P_{j,i}^{(m)}>0$.
\subsubsection{Equivalence Relation}
The concept of communication forms what is known as an equivalence relation, satisfying the following 
\begin{itemize}
    \item Reflexivity: $i\leftrightarrow i$
    \item Symmetry: $i\leftrightarrow j \Rightarrow j\leftrightarrow i$
    \item Transitivity: $i\leftrightarrow j$ and $j\leftrightarrow k \Rightarrow i\leftrightarrow k$
\end{itemize}
\subsubsection{Communication Class}
Within the same communication class, all states communicate with each other \\
If states $i$ and $j$ belongs to different communication classes, then $i\leftrightarrow j$ is not true, at most one of $i\to j$ or $j\to i$ can be true 
\subsubsection{Reducible, Irreducible}
A DTMC that has only one communication class is said to be irreducible. On the other hand, a DTMC is 
called reducible if there are two or more communication classes 
\subsubsection{Period}
The period of state $i$ is given by $d(i) = \gcd\{n\in\Z^+: P_{i,i}^{(n)}>0\}$
\subsubsection*{Aperiodic}
If $d(i)=1$, then state $i$ is said to be aperiodic. In fact, DTMC is said to be aperiodic if $d(i)=1$ for all $n\in\N$. Furthermore,
if $P_{i,i}^{(n)}=0$ for all $n\in\Z^+$, then we set $d(i)=\infty$ 
\subsubsection{Theorem}
If $i\leftrightarrow j$, then $d(i)=d(j)$
\subsection{Transience and Recurrence}
The first visit of the DTMC to state $j$ occurs at time $n\in\Z^+$ 
\begin{align*}
    f_{i,j}^{(n)} &= P(\text{DTMC makes a future visit to state } j \text{ at time } n|X_0=i) \\
    &= P(X_n=j, X_{n-1}\neq j, X_{n-2}\neq j, \cdots, X_1\neq j | X_0=i) \\
    f_{ii}^{(n)} &= P(X_n=i, X_{n-1}\neq i, X_{n-2}\neq i, \cdots, X_1\neq i | X_0=i) \\
    f_{ii}^{(0)} &= 0
\end{align*}
Result: 
\[f_{ij}^{(n)} = P_{ij}^{(n)} - \sum_{k=1}^{n-1}f_{ij}^{(k)}P_{jj}^{(n-k)}\]    
\subsubsection*{Definition}
\[f_{ij} = P(\text{DTMC ever makes a visit to state } j|X_0=i)\]    
\[T_i^{(n)}: \text{time for DRMC to makes n-th visit at state } i\]    
\[f_{i,j}^{(n)} = P(T_j^{(n)} = n | X_0=i)\]    
\subsubsection{Transient, Recurrent}
State $i$ is said to be transient if $f_{i,i}<1$. State $i$ is said to be recurrent if $f_{i,i}=1$
\subsubsection*{Definition}
$M_i$: number of times a DTMC make a visit to state $i$
\begin{align*}
    A_n^{(i)} = 
    \begin{cases}
        1 & \text{if } X_n=i \\
        0 & \text{if } X_n\neq i
    \end{cases}
\end{align*}
\[M_i = \sum_{n=0}^{\infty}A_n^{(i)}\]
Given $X_0=i$, 
\begin{itemize}
    \item if state $i$ is recurrent, then $M_i=\infty$ with probability $1$
    \item if state $i$ is transient, then $M_i\sim Geo_f(1-f_{ii})$
\end{itemize}
\subsubsection*{Definition}
$T_i^{(n)}$: time at n-th arrival to state $i$ 
\begin{align*} 
    I_i^{(n)} &= T_i^{(n)} - T_i^{(n-1)} \\
    &= \text{time between the } (n-1)^{th} \text{ and } n^{th} \text{ visit to state } i
\end{align*}
Define $T_i^{(0)}=0$ \\
Then $T_i^{(1)} = I_i^{(1)}$, $T_i^{(n)} = I_i^{(1)} + I_i^{(2)} + \cdots + I_i^{(n)}$ 
\[f_{ii} = P(T_i^{(1)}<\infty|X_0=i) = P(I_i^{(1)}<\infty|X_0=i)\]   
\subsubsection*{Theorem}
For a DTMC with stationary transition probability, the following holds
\begin{itemize}
    \item given $X_0=i$, the inter arrival time $I_1^{(1)}, I_1^{(2)},\cdots$ are iid integer valued random variables 
    \item if state $i$ is transient, then $M_i\sim Geo_f(1-f_{ii})$, where $M_i$ is number of arrivals at state $i$ 
    \item if state $i$ is recurrent, then $M_i = \infty$ with probability 1 
\end{itemize}
\subsubsection*{Corollary}
\begin{align*}
    M_i^{(n)} &= 
    \begin{cases}
        1 &\text{if } X_n=i \\
        0 &\text{if } X_n\neq i 
    \end{cases} \\
    &= \dso \{X_n=i\}
\end{align*}
\[M_i = \sum_{n=1}^{\infty}\dso \{X_n=i\}\]
Then 
\begin{align*}
    E[M_i|X_0=i] &= E\left[\sum_{n=1}^{\infty}\dso \{X_n=i\}|X_0=i\right] \\
                 &= \sum_{n=1}^{\infty}E\left[\dso \{X_n=i\}|X_0=i\right] \\
                 &= \sum_{n=1}^{\infty}P(X_n = i|X_0=i) \\
                 &= \sum_{n=1}^{\infty}P_{i,i}^{(n)}
\end{align*}
\begin{itemize}
    \item If state $i$ is transient \[E[M_i|X_0=i] = \sum_{n=1}^{\infty}P_{i,i}^{(n)} < \infty\]
    \item If state $i$ is recurrent \[E[M_i|X_0=i] = \sum_{n=1}^{\infty}P_{i,i}^{(n)} = \infty\]
    \item A state is recurrent iff \[\sum_{n=1}^{\infty}P_{i,i}^{(n)} = \infty\] 
    \item Symmetric random walk on $\Z$ is recurrent, but asymmetric random walk are transient 
\end{itemize}
\subsubsection*{Theorem}
If $i\leftrightarrow j$ and state $i$ is recurrent, then state $j$ is also recurrent \\
If $i\leftrightarrow j$ and state $i$ is transient, then state $j$ is also transient
\subsubsection*{Theorem}
If $i\leftrightarrow j$ and state $i$ is recurrent, then 
\[f_{i,j} = P(\text{DTMC ever makes a future visit to state } j|X_0=i)=1\] 
\subsubsection*{Finite DTMC}
DTMC is called finite if the total number of states is finite 
\subsubsection*{Theorem}
Consider a finite DTMC 
\begin{itemize}
    \item At least one of the states is recurrent 
    \item If the chain is irreducible, then the DTMC is recurrent, means all states of the DTMC are recurrent 
    \item If the chain is not irreducible, then it has at least one recurrence class 
\end{itemize}
A finite-state DTMC has at least one recurrent state \\
An irreducible, finite-state DTMC must be recurrent 
\subsubsection*{Theorem}
If state $i$ is recurrent and state $i$ does not communicate with state $j$, then $P_{i,j}^{(k)}=0$ for all $k\in\Z^+$ 
\subsubsection*{Theorem}
If $i\to j$ and $i$ is recurrent, then 
\begin{itemize}
    \item $j\to i$
    \item $i\leftrightarrow j$ and $j$ is recurrent 
    \item $f_{i,j} = P(X_n=j\text{ for some }n\geq 1|X_0=i) = 1$, $f_{j,i} = 1$
\end{itemize}
\subsection{Limiting Behaviour of DTMCs}
\subsubsection*{Theorem}
For any state $i$ and transient state $j$ of a DTMC, $\ds\lim_{n\to\infty}P_{i,j}^{(n)} = 0$
\subsubsection*{Definition}
Let 
\[N_i=\min\{n\in\Z^+:X_n=i\}\]   
where state $i$ is assumed to be recurrent. The first return time to state $i$ is defined as 
\[T_i^{(1)} = N_i\]   
The conditional pmf is given by 
\[P(N_i=n|X_0=i)=f_{i,i}^{(n)}, n=1,2,3,\cdots\]    
\subsubsection{Mean Recurrent Time}
If state $i$ is recurrent, then its mean recurrent time is given by 
\[m_i = E[N_i|X_0=i] = \sum_{n=1}^{\infty}n f_{i,i}^{(n)}\]   
In words, $m_i$ represents the average time it takes the DTMC to make successive visits to state $i$. 
\subsubsection{Positive Recurrent, Null Recurrent}
Suppose that state $i$ is recurrent. State $i$ is said to be positive recurrent if $m_i<\infty$. On the other hand, state $i$ is said to be null recurrent if $m_i=\infty$
\subsubsection*{Theorem}
\begin{itemize}
    \item If $i\leftrightarrow j$ and state $i$ is positive recurrent, then state $j$ is also positive recurrent 
    \item If $i\leftrightarrow j$ and state $i$ is null recurrent, then state $j$ is also null recurrent
    \item In a finite state DTMC, all recurrent states are positive recurrent
    \item In a finite state DTMC, at least one state is positive recurrent
\end{itemize}
\subsubsection{Stationary Distribution}
A probability distribution $\{p_i\}_{i=0}^{\infty}$ is called a stationary distribution of a DTMC if 
$\{p_i\}_{i=0}^{\infty}$ satisfies the conditions $\ds\sum_{i=0}^{\infty}p_i=1$ and $p_j=\ds\sum_{i=0}^{\infty}p_iP_{i,j}$ for all $j\in\N$
\[\pi P = \pi\]
\subsubsection{Convex Combination}
For any 2 vectors $\pi_1 = (\pi_{11},\pi_{12},\cdots,\pi_{1n})$ and $\pi_2 = (\pi_{21},\pi_{22},\cdots,\pi_{2n})$, a convex combination of $\pi_1$ and $\pi_2$ is given by
\[\alpha\pi_1 + (1-\alpha)\pi_2\quad 0\leq\alpha\leq 1\]
\subsubsection*{Result}
For any 2 stationary distributions $\pi_1$ and $\pi_2$, the convex combination $\alpha\pi_1 + (1-\alpha)\pi_2$ where $0\leq\alpha\leq1$ is another stationary distribution 
\begin{align*}
    (\alpha\pi_1 + (1-\alpha)\pi_2)P = \alpha\pi_1 + (1-\alpha)\pi_2
\end{align*}
\subsubsection*{Theorem}
$T_x^{(1)}$: time to first return to state $x$
\[m_x = E[T_x^{(1)}|X_0=x]\] 
\[f_{xy} = P(X_n = y, n\geq 1|X_0=x)\]   
\[\frac{1}{n}\sum_{n=1}^{n}P_{x,y}^{(n)}\rightarrow \frac{f_{xy}}{m_y}\quad n\to\infty\]   
\begin{itemize}
    \item $m_y = \infty$ if $y$ is transient or null recurrent
    \item $m_y < \infty$ if $y$ is positive recurrent
\end{itemize}
\subsubsection*{Theorem}
If the DTMC is 
\begin{itemize}
    \item irreducible 
    \item aperiodic 
    \item positive recurrent 
\end{itemize}
Then the stationary distribution is unique and 
\[\lim_{n\to\infty}P_{x,y}^{(n)} = \pi_y = \frac{1}{m_y}\]    
\[\lim_{n\to\infty}P(X_n=x) = \pi_x = \frac{1}{m_x}\]    
For a positive recurrent state $y$, the expected return time is 
\[m_y = \frac{1}{\pi_y}\]    
\subsubsection*{Result}
\begin{itemize}
    \item Stationary distribution exists iff at least one of the state is positive recurrent
    \begin{itemize}
        \item For finite DTMC, at least one state is positive recurrent, hence stationary distribution always exists 
    \end{itemize}
    \item If more than one positive recurrent classes, then stationary distribution is not unique 
    \item For finite chain the stationary distribution is unique if it has only one positive recurrent class 
\end{itemize}
\subsubsection{Basic Limit Theorem}
For an irreducible, recurrent, and aperiodic DTMC, $\ds\lim_{n\to\infty}P_{i,j}^{(n)}$ exists 
and is independent of state $i$, satisfying 
\[\lim_{n\to\infty}P_{i,j}^{(n)}=\pi_j=\frac{1}{m_j}\quad\forall i,j\in\N\]    
If the DTMC also happens to be positive recurrent, then $\{\pi_j\}_{j=0}^{\infty}$ is the unique, positive solution
to the system of linear equations defined by 
\begin{align*}
    \begin{cases}
        \pi_j = \sum_{i=0}^{\infty}\pi_i P_{i,j} & \forall j\in\N \\
        \sum_{j=0}^{\infty}\pi_j = 1
    \end{cases}
\end{align*}
\subsubsection*{Theorem}
Suppose that a finite-state DTMC with state space $S=\{0,1,\cdots,N-1\}$ is irreducible and aperiodic.
If the associated TPM is doubly stochastic, then the limiting probabilities $\{\pi_j\}_{j=0}^{N-1}$ exist 
and are given by 
\[\pi_j = \frac{1}{N}\quad j=0,1,\cdots,N-1\]    
\subsubsection{Probability Generating Function}
\begin{align*}
    \alpha(t) &= \sum_{j=0}^{\infty}t^j\alpha_j \\
    &= E[t^{Z_{0,1}}]
\end{align*}
\[\frac{d^n}{dt^n}(\alpha(t))\bigg|_{t=0} = (n!)\cdot \alpha_n\]
Let $\pi_0 = P(X_n=0\text{ for some } n\geq1|X_0=1)$ 
\begin{itemize}
    \item $1$ is a solution to $\pi_0=\alpha(\pi_0)$
    \item $\pi_0$ is a solution to the equation \[t=\alpha(t)\]
    \item $f(t) = \alpha(t)-t$ is a convex function in $[0,1]$
    \item Since $f(t) = \alpha(t)-t$ is strictly convex, smooth and non constant inside $[0,1]$, the equation $f(t)=0$ has at most two solution in $[0,1]$
\end{itemize}
\subsubsection*{Theorem}
$\pi_0 = \ds\lim_{n\to\infty}P(X_n=0)$ is the smallest solution of $t=\alpha(t)$ within the interval $[0,1]$

\newpage 

\section{The Exponential Distribution and the Poisson Process}
\subsection{Properties of the Exponential Distribution}
If a rv $X$ has an exponential distribution with parameter $\lambda>0$ ($X\sim Exp(\lambda)$ where $\lambda$ is often referred to as the "rate"),
then we have the following distributional results 
\begin{itemize}
    \item pdf: $f(x) = \lambda e^{-\lambda x}, x\geq 0$
    \item cdf: $F(x) = 1-e^{-\lambda x}, x\geq 0$
    \item tpf: $\overline{F}(x) = e^{-\lambda x}, x\geq 0$
    \item mgf: $\phi_X(t) = \frac{\lambda}{\lambda-t}, t<\lambda$
    \item mean: $E[X] = \frac{1}{\lambda}$
    \item variance: $Var(X) = \frac{1}{\lambda^2}$
\end{itemize}
\subsubsection*{Minimum of Independent Exponential}
Let $\{X_i\}_{i=1}^n$ be a sequence of independent rvs where $X_i\sim Exp(\lambda_i)$, $i=1,2,\cdots,n$. Define $Y=\min\{X_1,X_2,\cdots,X_n\}$ to be the smallest order statistic of $\{X_1,\cdots,X_n\}$. Then 
\[Y=\min\{X_1,X_2,\cdots,X_n\}\sim Exp(\sum_{i=1}^{n}\lambda_i)\]
\subsubsection*{Usuful Result}
Let $\{X_i\}_{i=1}^n$ be a sequence of independent rvs where $X_i\sim Exp(\lambda_i)$, $i=1,2,\cdots,n$ 
\begin{itemize}
    \item For $j\in\{1,2,\cdots,n\}$ \[P(X_j = \min\{X_1,X_2,\cdots,X_n\})  = \frac{\lambda_j}{\sum_{i=1}^{n}\lambda_i}\]
    \item \[P(X_1<X_2<\cdots<X_n) = \frac{\prod_{i=1}^{n-1}\lambda_i}{\prod_{i=1}^{n-1}\left(\sum_{j=i}^{n}\lambda_j\right)}\]
\end{itemize}
\subsubsection{Memoryless Property}
A rv $X$ is memoryless iff 
\[P(X>y+z|X>y) = P(X>z)\quad\forall y,z\geq 0\]
$X$ is memoryless iff 
\[P(X>y+z) = P(X>y)P(X>z)\quad\forall y,z\geq 0\]
\subsubsection*{Theorem}
An exponential distribution is memoryless
\subsubsection*{Theorem}
If $X_1\sim Exp(\lambda_1)$ and $X_2\sim Exp(\lambda_2)$ are independent rvs, then
\[P(X_1<X_2) = \frac{\lambda_1}{\lambda_1+\lambda_2}\] 
\subsubsection*{Useful Results}
If $X\sim Exp(\lambda)$
\[P(X>Y+Z|X>Y) = P(X>Z)\]
If $X\sim Exp(\lambda)$, then 
\[P(X-Y>z|X>Y) = P(X>z)\]
\[(X-Y)|(X>Y)\sim Exp(\lambda)\]
\subsection{The Poisson Process}
A counting process $\{N(t),t\geq0\}$ is a stochastic process in which $N(t)$ represents the number of events that happen (or occur) by time $t$, where the index $t$ measures time over a continuous range 
\subsubsection*{Basic Properties of Counting Processes}
\begin{itemize}
    \item $N(0)=0$
    \item $N(t)$ is a non-negative integer $\forall t\geq0$ ($N(t)\in\N\quad\forall t\geq0$)
    \item If $s<t$, then $N(s)\leq N(t)$
    \item $N(t)-N(s)$ counts the number of events to occur in the time interval $(s,t]$ for $s<t$
\end{itemize}
\subsubsection{Independent Increments}
A counting process $\{N(t),t\geq0\}$ has independent increments if $N(t_1)-N(s_1)$ is independent of 
$N(t_2)-N(s_2)$ whenever $(s_1,t_1]\cap(s_2,t_2]=\emptyset$ for all choices of $s_1,t_1,s_2,t_2$
\subsubsection{Stationary Increments}
A counting process $\{N(t),t\geq0\}$ has stationary increments if the distribution of the number of 
events in $(s,s+t]$ depends only on $t$, the length of the time interval. $N(s+t)-N(s)$ has the same probability distribution as $N(t)-N(0)=N(t)$, the number of events occurring in the interval $[0,t]$
\subsubsection{Of order $h$}
A function $y=f(x)$ is said to be "$o(h)$" (of order $h$) if 
\[\lim_{h\to0} \frac{f(x)}{h} = 0\]
\subsubsection{Poisson Process}
A counting process $\{N(t),t\geq 0\}$ is said to be a Poisson process at rate $\lambda$ if the following three conditions hold true 
\begin{itemize}
    \item The process has both independent and stationary increments 
    \item For $h>0$, $P(N(h)=1) = \lambda h + o(h)$
    \item For $h>0$, $P(N(h)\geq 2)=o(h)$
\end{itemize}
\subsubsection*{Theorem}
If $\{N(t),t\geq0\}$ is a Poisson process at rate $\lambda$, then 
\[N(t)\sim Poi(\lambda t)\]
\subsubsection{Interarrival Times}
Define $T_1$ to be the elapsed time (from time $0$) until the first event occurs. In general, for $i\geq2$, let $T_i$ be the elapsed time between 
the occurences of the $(i-1)^{th}$ and $i^{th}$ events. The sequence $\{T_i\}_{i=1}^{\infty}$ is called
the interarrival or interevent time sequence 
\subsubsection*{Theorem}
If $\{N(t),t\geq0\}$ is a Poisson process at rate $\lambda>0$, then $\{T_i\}_{i=1}^{\infty}$ is a sequence 
of i.i.d. $Exp(\lambda)$ rvs 
\[S_n = \sum_{i=1}^{n}T_i\sim Erlang(n,\lambda)\]
\[P(S_n>t) = e^{-\lambda t}\sum_{j=0}^{n-1}\frac{(\lambda t)^j}{j!}, t\geq 0\]
\subsection{Further Properties of the Poisson Process}
\subsubsection*{Theorem}
If $\{N(t),t\geq0\}$ is a Poisson process at rate $\lambda$, then 
\[N(s)|(N(t)=n)\sim Bin(n,s/t)\]
where $s<t$
\subsubsection{Comparison of Event Occurences}
Suppose that $\{N_1(t),t\geq0\}$ and $\{N_2(t),t\geq0\}$ are independent Poisson process at rate $\lambda_1$ and $\lambda_2$ respectively. 
Let $S_m^{(1)}$ be the arrival time of the $m$-th event for $\{N_1(t),t\geq0\}$ and $S_n^{(2)}$ be the arrival time of the $n$-th event for $\{N_2(t),t\geq0\}$. 
We know that $S_m^{(1)} = \sum_{i=1}^{m}T_i^{(1)}$ where $\{T_i^{(1)}\}_{i=1}^{\infty}$ is a sequence of iid $Exp(\lambda_1)$ rvs and $S_n^{(2)} = \sum_{j=1}^{n}T_j^{(2)}$ where $\{T_j^{(2)}\}_{j=1}^{\infty}$ is a sequence of iid $Exp(\lambda_2)$ rvs. 
Then the following holds true\\
The probability of observing $m$ success occur before $n$ failure in a sequence of independent Bernoulli trials
\[P(S_m^{(1)}<S_n^{(2)}) = \sum_{j=0}^{n-1}\binom{m+j-1}{m-1}\left(\frac{\lambda_1}{\lambda_1+\lambda_2}\right)^m\left(\frac{\lambda_2}{\lambda_1+\lambda_2}\right)^j\]
\subsubsection*{Theorem}
Suppose that $\{N(t),t\geq0\}$ is a Poisson process at rate $\lambda$. Given $N(t)=1$, the conditional 
distributino of the first arrival time is uniformly distributed on $(0,t)$ 
\[S_1|(N(t)=1)\sim U(0,t)\]
\subsubsection{Order Statistics}
The sequence of rvs $\{Y_{(1)},Y_{(2)},\cdots,Y_{(n)}\}$ are called the order statistics of $\{Y_1,Y_2,\cdots,Y_n\}$, satisfying 
\begin{align*}
    Y_{(1)} &\equiv \text{1st smallest among } \{Y_1,Y_2,\cdots,Y_n\} \\
    Y_{(2)} &\equiv \text{2nd smallest among } \{Y_1,Y_2,\cdots,Y_n\} \\
    &\vdots \\
    Y_{(n)} &\equiv \text{nth smallest among } \{Y_1,Y_2,\cdots,Y_n\}
\end{align*}
\subsubsection*{Theorem}
Let $\{N(t),t\geq0\}$ be a Poisson process at rate $\lambda$. Given $N(t)=n$, the conditional joint distribution 
of the $n$ arrival times is identical to the joint distribution of the $n$ order statistics from the $U(0,t)$ distributino 
\[(S_1,S_2,\cdots,S_n)|(N(t)=n)\sim (Y_{(1)},Y_{(2)},\cdots,Y_{(n)})\]
where $\{Y_i\}_{i=1}^{n}$ is an i.i.d. sequence of $U(0,t)$ rvs
\subsection{Two Important Generalizations}
\subsubsection{The Non-Homogeneous Poisson Process}
The counting process $\{N(t),t\geq0\}$ is a non-homogeneous (non-stationary) Poisson process with rate function $\lambda(t)$ if the 
following three conditions hold true 
\begin{itemize}
    \item $\{N(t),t\geq0\}$ has independent increments
    \item $P(N(t+h)-N(t)=1) = h\lambda(t) + o(h)$
    \item $P(N(t+h)-N(t)\geq 2)=o(h)$
\end{itemize}
\subsubsection*{Theorem}
If $\{N(t),t\geq0\}$ is a non-homogeneous Poisson process with rate functionn $\lambda(t)$, then 
$N(s_1+s_2)-N(s_1)\sim Poi(m(s_1+s_2)-m(s_1))$ where the mean value function $m(t)$ is given by 
\[m(t) = \int_{0}^{t}\lambda(\tau)d\tau,\quad t\geq0\]
\subsubsection{The Compound Poisson Process}
Let $\{Y_i\}_{i=1}^{\infty}$ be an i.i.d. sequence of rvs. Let $\{N(t),t\geq0\}$ be a Poisson process 
at rate $\lambda$, independent of each $Y_i$, $i=1,2,\cdots$. If $X(t)=\sum_{i=1}^{N(t)}Y_i$, then the 
process $\{X(t),t\geq0\}$ is a compound Poisson process 

\[E[X(t)] = \lambda tE(Y_1)\]
\[Var(X(t)) = \lambda tE(Y_1^2)\]

\newpage 
\section{Useful References}
\subsection{Ratio Test}
For a sequence of positive number $\{a_n\}$
\begin{itemize}
    \item If $L = \ds\lim_{n\to\infty}\frac{a_{n+1}}{a_n} < 1$, then $\sum_{n=1}^{\infty}a_n < \infty$ 
    \item If $L > 1$, then $\sum_{n=1}^{\infty}a_n = \infty$
    \item If $L = 1$, the test is inconclusive
\end{itemize}
\subsection{Stirling's Approximation}
\[n! \sim \sqrt{2\pi}e^{-n} n^{n+\frac{1}{2}}\]



\end{document}
